\documentclass[a4paper,12pt]{ltjsbook}
\usepackage{subfiles}
\usepackage{omebook}
\usepackage{svg}
\svgsetup{
  inkscapelatex=false
}
\setcounter{chapter}{6}
\setcounter{secnumdepth}{4}
\begin{document}
  \input body.tex
  \begin{thebibliography}{99}
  \bibitem{hsp} おにたま, 悠黒喧史, うすあじ. HSPでつくる! はじめてのプログラミング HSP3.5+3Dish入門. 秀和システム. 2017
  \\ OpenJtalkについて
  \bibitem{jtalk1} Open JTalk - HMM-based Text-to-Speech System \\ \url{http://open-jtalk.sp.nitech.ac.jp/}
  \bibitem{jtalk2} Open JTalk (テキストを音声へ\ruby{変換}{へん|かん})をインストールし、使ってみる | レンタルサーバー・自\ruby{宅}{たく}サーバー\ruby{設定}{せっ|てい}・\ruby{構築}{こう|ちく}のヒント \\ \url{https://server-setting.info/centos/open-jtalk-install.html}
  \\ Juliusについて
\bibitem{julius} 大語\ruby{彙}{い}連続音声\ruby{認識}{にん|しき}エンジン Julius \\ \url{http://julius.osdn.jp/}
  \end{thebibliography}
\end{document}
